\chapter{Desarrollo de la biblioteca wtActuatorFoam}

\section{Estado del arte}

\section{Parámetros generales de funcionamiento}

\begin{itemize}
	\item Mallado local del actuador
	\item Modelo de inducción
	\item Factor de corrección de punta de pala
	\item Factor de corrección de raíz de pala
	\item Cálculo de velocidades y distribución de fuerzas
	\item Niveles de guardado
\end{itemize}

\section{Modelos de actuador implementados}

\subsection{\texttt{uniform}}
\subsection{\texttt{airfoil}}
\subsection{\texttt{numeric}}
\subsection{\texttt{adaptive}}
\subsection{\texttt{analytic}}
\subsection{\texttt{genalytic}}

\section{Análisis de Escalabilidad}

\section{Verificación}

Acá vamos a hablar de las validaciones realizadas con otros modelos de actuadores presentados previamente, todo
lo que fue presentado en WESC2025 y en TORQUE2026.

\section{Caso de Aplicación: Análisis de rotación de la estela lejana en contexto de ASL}

\section{Conclusiones del Capítulo}


